\appendix
\part{Anhang}

\newcommand{\lit}[3]{\vspace{5pt}\noindent\textbf{[#1]} #2\par\nopagebreak\vspace{1pt}{\small\itshape\color{grau} #3}\par\vspace{2pt}}

\chapter{Formelsammlung und Übersicht}

\section{Beschreiben}

\begin{merksatz}
\boxtitel{Lage und Streuung}
\begin{center}
{\small
\renewcommand{\arraystretch}{1.6}
\begin{tabular}{@{}P{4.4cm}P{4.2cm}@{}}
\toprule
$\MW = \frac{1}{n}\sum_i x_i$ & arithmetisches Mittel (2.2)\\
$\sum_i (x_i-\MW) = 0$ & Satz 2.3\\
$\sum(x_i-a)^2$ minimal bei $a=\MW$ & Satz 2.4\\
$\sum|x_i-a|$ minimal bei $a=\Med$ & Satz 2.6\\
$s^2 = \frac{1}{n-1}\sum (x_i-\MW)^2$ & Stichprobenvarianz\\
$\sum(x_i-\MW)^2 = \sum x_i^2 - n\MW^2$ & Verschiebungssatz (2.11)\\
$\mathrm{IQR} = Q_3 - Q_1$ & Quartilsabstand\\
$v = s/\MW$ & Variationskoeffizient\\
$z_i = (x_i-\MW)/s$ & Standardisierung\\
\bottomrule
\end{tabular}}
\end{center}
\end{merksatz}

\begin{merksatz}
\boxtitel{Zusammenhang}
\begin{center}
{\small
\renewcommand{\arraystretch}{1.6}
\begin{tabular}{@{}P{4.4cm}P{4.2cm}@{}}
\toprule
$s_{xy} = \frac{1}{n-1}\sum(x_i-\MW)(y_i-\overline y)$ & Kovarianz (3.1)\\
$\sum(x-\MW)(y-\overline y) = \sum xy - n\MW\overline y$ & Satz 3.2\\
$r = s_{xy}/(s_xs_y)$ & Korrelation (3.3)\\
$-1 \le r \le 1$ & Satz 3.4\\
$r_s = 1 - \frac{6\sum d_i^2}{n(n^2-1)}$ & Spearman, ohne Bindungen\\
\bottomrule
\end{tabular}}
\end{center}
\end{merksatz}

\section{Schätzen und Testen}

\begin{merksatz}
\boxtitel{Grundlagen}
\begin{center}
{\small
\renewcommand{\arraystretch}{1.6}
\begin{tabular}{@{}P{4.4cm}P{4.2cm}@{}}
\toprule
$\mathbb{E}[\overline X] = \mu$ & Satz 4.4\\
$\Var[\overline X] = \sigma^2/n$ & Satz 4.4, braucht Unabhängigkeit\\
$\mathrm{SE} = s/\sqrt n$ & Standardfehler (4.5)\\
$\mathbb{E}[X^2] = \Var[X] + \mu^2$ & Hilfssatz zu 5.3\\
$\mathbb{E}[S^2] = \sigma^2$ & Satz 5.3 --- das $n-1$\\
$\mathrm{MSE} = \Var + \mathrm{Bias}^2$ & Satz 5.6\\
$\hat\pi = k/n$ & ML für den Anteil\\
$\hat\mu = \MW$, $\hat\sigma^2 = \frac1n\sum(x-\MW)^2$ & ML, normal --- verzerrt\\
$\hat\lambda = 1/\MW$ & ML, exponentiell\\
\bottomrule
\end{tabular}}
\end{center}
\end{merksatz}

\begin{merksatz}
\boxtitel{Intervalle und Tests}
\begin{center}
{\small
\renewcommand{\arraystretch}{1.6}
\begin{tabular}{@{}P{4.4cm}P{4.2cm}@{}}
\toprule
$\MW \pm t_{n-1}\,\frac{s}{\sqrt n}$ & $t$-Intervall (6.2)\\
$s_p^2 = \frac{(n_A-1)s_A^2+(n_B-1)s_B^2}{n_A+n_B-2}$ & gepoolte Varianz\\
$\mathrm{SE} = s_p\sqrt{\frac{1}{n_A}+\frac{1}{n_B}}$ & Student\\
$\mathrm{SE} = \sqrt{\frac{s_A^2}{n_A}+\frac{s_B^2}{n_B}}$ & Welch\\
Wilson & statt Wald bei Anteilen\\
$3/n$ & Dreierregel bei null Ereignissen\\
$n = \frac{2(z_{1-\alpha/2}+z_{1-\beta})^2\sigma^2}{\delta^2}$ & Fallzahl je Gruppe\\
$\chi^2 = \sum \frac{(O-E)^2}{E}$ & Vierfeldertafel\\
$U_A = R_A - \frac{n_A(n_A+1)}{2}$ & Mann-Whitney\\
\bottomrule
\end{tabular}}
\end{center}
\end{merksatz}

\section{Modelle und Medizin}

\begin{merksatz}
\boxtitel{Regression}
\begin{center}
{\small
\renewcommand{\arraystretch}{1.6}
\begin{tabular}{@{}P{4.4cm}P{4.2cm}@{}}
\toprule
$b = \frac{\sum(x-\MW)(y-\overline y)}{\sum(x-\MW)^2}$ & Satz 10.2\\
$a = \overline y - b\MW$ & Gerade durch den Schwerpunkt\\
$b = r\,s_y/s_x$ & Zusammenhang mit $r$\\
$R^2 = r^2$ & nur bei einfacher Regression\\
$\widehat{\mathrm{SE}}(b) = s_e/\sqrt{\sum(x-\MW)^2}$ & mit $s_e^2 = SS_{\text{Rest}}/(n-2)$\\
$\hat z_y = r\,z_x$ & Regression zur Mitte\\
\bottomrule
\end{tabular}}
\end{center}
\end{merksatz}

\begin{merksatz}
\boxtitel{Andere Zielgrößen und Diagnostik}
\begin{center}
{\small
\renewcommand{\arraystretch}{1.6}
\begin{tabular}{@{}P{4.4cm}P{4.2cm}@{}}
\toprule
$\operatorname{logit}(\pi) = \ln\frac{\pi}{1-\pi}$ & logistisches Modell (11.2)\\
$e^{\beta}$ & Chancenverhältnis je Einheit\\
$\mathrm{OR} = \frac{ad}{bc}$ & Kreuzprodukt\\
$\widehat{\mathrm{SE}}(\ln\mathrm{OR}) = \sqrt{\frac1a+\frac1b+\frac1c+\frac1d}$ & \\
$S(t) = \prod_{t_j \le t}(1-d_j/n_j)$ & Kaplan-Meier\\
$h(t\mid x) = h_0(t)e^{\beta x}$ & Cox (11.3)\\
$\mathrm{LQ}^{+} = \frac{\text{Sens}}{1-\text{Spez}}$ & Likelihood-Quotient\\
Chance nachher $=$ Chance vorher $\times\ \mathrm{LQ}$ & Bayes, Chancenform\\
$\mathrm{AUC} = U/(n_1n_2)$ & wie Mann-Whitney\\
\bottomrule
\end{tabular}}
\end{center}
\end{merksatz}

\begin{merksatz}
\boxtitel{Viele Tests, viele Studien}
\begin{center}
{\small
\renewcommand{\arraystretch}{1.6}
\begin{tabular}{@{}P{4.4cm}P{4.2cm}@{}}
\toprule
$1-(1-\alpha)^m$ & Fehlerrate bei $m$ Tests\\
$\alpha/m$ & Bonferroni\\
$\alpha/(m-i+1)$ & Holm, aufsteigend\\
$\frac{i}{m}q$ & Benjamini-Hochberg\\
$w_i = 1/\mathrm{SE}_i^2$ & Metaanalyse, fester Effekt\\
$Q = \sum w_i(\hat\theta_i - \hat\theta)^2$ & Heterogenität\\
$I^2 = \max(0, \frac{Q-(k-1)}{Q})$ & \\
$w_i^{*} = 1/(\mathrm{SE}_i^2+\tau^2)$ & zufälliger Effekt\\
\bottomrule
\end{tabular}}
\end{center}
\end{merksatz}

\section{Die Landkarte des Bandes}

\begin{merksatz}
\boxtitel{Welcher Satz wo --- und worauf er beruht}
\begin{center}
{\small
\renewcommand{\arraystretch}{1.42}
\begin{tabular}{@{}P{3.4cm}P{1.2cm}P{3.6cm}@{}}
\toprule
\textbf{Ergebnis} & \textbf{Kap.} & \textbf{Beruht auf}\\
\midrule
Skalenniveaus & $1$ & erlaubte Umformungen\\
Mittelwert minimiert Quadrate & $2$ & Null einfügen\\
Median minimiert Beträge & $2$ & Verschiebeargument\\
Verschiebungssatz & $2$ & Ausquadrieren\\
$-1 \le r \le 1$ & $3$ & Summe von Quadraten\\
Simpson & $3$ & gewichtetes Mittel\\
$\Var[\overline X] = \sigma^2/n$ & $4$ & Unabhängigkeit\\
Grenzwertsatz & $4$ & \textcolor{grau}{nicht bewiesen}\\
$\mathbb{E}[S^2] = \sigma^2$ & $5$ & Satz 4.4 auf $\overline X$\\
MSE-Zerlegung & $5$ & derselbe Kniff wie 2.4\\
ML für $\mu$ & $5$ & \textbf{Satz 2.4}\\
$t$-Intervall & $6$ & Umformung der Ungleichung\\
Überlappungs\-problem & $6$ & $\sqrt{a^2+b^2} < a+b$\\
Sieben Fehldeutungen & $7$ & $\Pr[D|H_0]$ gegen $\Pr[H_0|D]$\\
Varianzanalyse & $8$ & Zerlegung wie 2.4\\
Fisher exakt & $8$ & Abzählen\\
Permutationstest & $9$ & Austauschbarkeit\\
Kleinste Quadrate & $10$ & Ableiten, Satz 2.3\\
Regression zur Mitte & $10$ & $b = rs_y/s_x$\\
Chancenverhältnis & $11$ & Logit-Differenz\\
Kaplan-Meier & $11$ & Produkt bedingter Anteile\\
Fall-Kontroll & $12$ & Spaltenweise Bildung\\
Prävalenz\-abhängigkeit & $13$ & Satz von Bayes\\
AUC $= U/(n_1n_2)$ & $13$ & Kapitel~9\\
Bonferroni & $14$ & Vereinigungs\-abschätzung\\
\bottomrule
\end{tabular}}
\end{center}
\end{merksatz}

\begin{merksatz}
\boxtitel{Was in diesem Band \kursiv{nicht} bewiesen wurde}
\begin{center}
{\small
\renewcommand{\arraystretch}{1.5}
\begin{tabular}{@{}P{3.4cm}P{1.4cm}P{3.4cm}@{}}
\toprule
\textbf{Ergebnis} & \textbf{Kap.} & \textbf{Wo es steht}\\
\midrule
Zentraler Grenzwertsatz & $4$ & Band 4, Teil IV\\
Form der $t$-Verteilung & $6$ & Casella/Berger\zit{7}\\
Zweite Momentenmethode & $9$ & Efron/Tibshirani\zit{11}\\
Partielle Likelihood & $11$ & Cox\zit{9}\\
Tschebyschew-Ungleichung & $5$ & Band 4, Kapitel 7\\
\bottomrule
\end{tabular}}
\end{center}

\begin{center}
\begin{minipage}{0.88\textwidth}
\kursiv{Fünf Auslassungen. Jede ist im Text als solche benannt, mit Begründung und Fundstelle. Das ist die einzige zulässige Form, die Regel dieses Werkes zu brechen: offen und mit Quellenangabe.}
\end{minipage}
\end{center}
\end{merksatz}

\section{Zeittafel}

\begin{merksatz}
\boxtitel{Von den Zitronen bis zur Replikationskrise}
\begin{center}
{\small
\renewcommand{\arraystretch}{1.36}
\begin{tabular}{@{}P{1.4cm}P{2.8cm}P{4.0cm}@{}}
\toprule
\textbf{Jahr} & \textbf{Person} & \textbf{Ereignis}\\
\midrule
$1747$ & Lind & erster kontrollierter Versuch\\
$1763$ & Bayes & Satz über umgekehrte Wahrscheinlichkeit\\
$1774$ & Laplace & allgemeine Form\\
$1805$ & Legendre & kleinste Quadrate\\
$1809$ & Gauß & Normalverteilung als Fehlergesetz\\
$1835$ & Quetelet & der Durchschnittsmensch\\
$1858$ & Nightingale & Statistik als Argument\\
$1886$ & Galton & Regression zur Mitte\\
$1896$ & Pearson & Korrelationskoeffizient\\
$1900$ & Pearson & $\chi^2$-Test\\
$1908$ & Gosset & $t$-Verteilung\\
$1922$ & Fisher & Maximum-Likelihood\\
$1925$ & Fisher & Signifikanztest, Varianzanalyse\\
$1928$ & Neyman, Pearson & Entscheidungsverfahren\\
$1935$ & Fisher & \kursiv{Design of Experiments}\\
$1937$ & Neyman & Konfidenzintervall\\
$1945$ & Wilcoxon & Rangtests\\
$1946$ & Stevens & Skalenniveaus\\
$1947$ & Mann, Whitney & $U$-Test\\
$1948$ & Bradford Hill & erste randomisierte Studie\\
$1956$ & Stein & Zusammenziehen\\
$1958$ & Kaplan, Meier & Überlebenskurve\\
$1965$ & Bradford Hill & neun Kriterien\\
$1972$ & Cox & Regressionsmodell\\
$1972$ & Cochrane & \kursiv{Effectiveness and Efficiency}\\
$1976$ & Glass & Metaanalyse\\
$1977$ & Tukey & \kursiv{Exploratory Data Analysis}\\
$1979$ & Efron & Bootstrap\\
$1993$ & --- & Cochrane-Zusammenarbeit\\
$1995$ & Benjamini, Hochberg & Falscherkennungsrate\\
$2001$ & Breiman & die zwei Kulturen\\
$2005$ & Ioannidis & \kursiv{Most findings are false}\\
$2005$ & --- & Registrierungspflicht\\
$2009$ & Pearl & \kursiv{Causality}, 2.\ Auflage\\
$2015$ & Open Science & Wiederholungsvorhaben\\
$2016$ & ASA & Erklärung zum $p$-Wert\\
$2019$ & Amrhein u.\,a. & \enquote{signifikant} abschaffen\\
\bottomrule
\end{tabular}}
\end{center}
\end{merksatz}

\section{Glossar}

\begin{merksatz}
\boxtitel{Begriffe von A bis Z}
\begin{center}
{\small
\renewcommand{\arraystretch}{1.3}
\begin{longtable}{@{}P{3.2cm}P{5.4cm}@{}}
\toprule
\textbf{Begriff} & \textbf{Bedeutung}\\
\midrule
\endhead
Alpha & vorab festgelegtes Niveau (7.3)\\
AUC & Fläche unter der ROC-Kurve (13.2)\\
Ausreißer & auffällig abweichender Wert (2.4)\\
Austauschbarkeit & Voraussetzung des Permutationstests (9)\\
Bias & systematische Abweichung (5.2)\\
Bootstrap & Ziehen aus der Stichprobe (9)\\
Bruchpunkt & Robustheitsmaß (2.9)\\
Chance & $\pi/(1-\pi)$ (11.1)\\
Erwartungstreu & $\mathbb{E}[\hat\theta] = \theta$ (5.2)\\
Fallzahl & Umfang je Gruppe (7)\\
Freiheitsgrade & Zahl unabhängiger Abweichungen\\
Hazard & momentane Ereignisrate (11.3)\\
Heterogenität & echte Verschiedenheit der Studien (14)\\
$I^2$ & Anteil nicht zufälliger Streuung (14)\\
Intention-to-treat & Auswertung wie randomisiert (12)\\
Kollider & von beiden beeinflusst (3, 15)\\
Konfidenzintervall & Überdeckung des Verfahrens (6.1)\\
Konfounder & wirkt auf Ursache und Wirkung (3)\\
Konsistenz & Annäherung bei wachsendem $n$ (5.4)\\
Kovarianz & ungenormtes Zusammenhangsmaß (3.1)\\
Likelihood & Wahrscheinlichkeit der Daten als Funktion von $\theta$ (5)\\
Macht & $1-\beta$ (7.3)\\
Median & mittlere Ausprägung (2.5)\\
Mediator & liegt auf dem Wirkungsweg (3)\\
Metaanalyse & Zusammenfassung mehrerer Studien (14)\\
Multiplizität & viele Tests zugleich (14)\\
Nullhypothese & Annahme ohne Effekt (7.1)\\
Ordinalskala & geordnet, Abstände undeutbar (1.2)\\
$p$-Wert & $\Pr[\text{so extrem} \mid H_0]$ (7.2)\\
Perzentil & siehe Quantil (2.7)\\
Prädiktiver Wert & bedingt auf den Befund (13.1)\\
Prävalenz & Häufigkeit in der Grundgesamtheit (13)\\
Pseudoreplikation & abhängige Werte als unabhängig (4)\\
Randomisierung & zufällige Zuteilung (12)\\
Residuum & $y_i - \hat y_i$ (10.3)\\
Robustheit & Unempfindlichkeit gegen Ausreißer (2.4)\\
ROC & Sensitivität gegen $1-$Spezifität (13)\\
Sensitivität & bedingt auf Krankheit (13.1)\\
Signifikanz & $p \le \alpha$ --- nicht Bedeutsamkeit (7)\\
Spezifität & bedingt auf Gesundheit (13.1)\\
Standardfehler & Streuung des Schätzers (4.5)\\
Stichprobe & untersuchte Teilmenge (4.1)\\
$\tau^2$ & Zusatzvarianz zwischen Studien (14)\\
Verblindung & Unkenntnis der Zuteilung (12)\\
Verteilungsfrei & ohne Normalannahme (9)\\
Zensierung & Beobachtung endete vorzeitig (11)\\
$z$-Wert & standardisierter Wert (2.12)\\
\bottomrule
\end{longtable}}
\end{center}
\end{merksatz}
